\documentclass[12pt,a4paper]{article}

% Packages
\usepackage[utf8]{inputenc}
\usepackage{amsmath, amssymb}
\usepackage{graphicx}
\usepackage{booktabs}
\usepackage{hyperref}

% Title information
\title{A Simple \LaTeX{} Test Document}
\author{John Doe}
\date{\today}

\begin{document}

\maketitle

\begin{abstract}
This document is a simple test file for verifying that a \LaTeX{} environment
is correctly installed and functioning. It demonstrates basic structures
commonly used in academic writing.
\end{abstract}

\section{Introduction}

\LaTeX{} is a typesetting system widely used for scientific and technical
documents. It is especially popular in mathematics, computer science,
and engineering due to its powerful handling of formulas and references.

\section{Mathematical Example}

An example of an inline equation is $E = mc^2$.
A displayed equation is shown below:
\begin{equation}
\label{eq:sample}
\int_{0}^{\infty} e^{-x^2} \, dx = \frac{\sqrt{\pi}}{2}.
\end{equation}

Equation~\ref{eq:sample} is a well-known result in mathematics.

\section{Lists}

An example of an itemized list:
\begin{itemize}
  \item First item
  \item Second item
  \item Third item
\end{itemize}

An example of an enumerated list:
\begin{enumerate}
  \item Step one
  \item Step two
  \item Step three
\end{enumerate}

\section{Table}

Table~\ref{tab:example} shows a simple table.

\begin{table}[h]
\centering
\begin{tabular}{lcc}
\toprule
Method & Accuracy & Time (s) \\
\midrule
Method A & 0.85 & 1.2 \\
Method B & 0.89 & 1.8 \\
Method C & 0.83 & 0.9 \\
\bottomrule
\end{tabular}
\caption{An example table for testing purposes}
\label{tab:example}
\end{table}

\section{Conclusion}

This test document confirms that text, mathematics, lists, tables,
and references can be compiled successfully. It can be extended
as a starting point for more complex academic documents.

\bibliographystyle{plain}
\begin{thebibliography}{1}
\bibitem{lamport1994}
Leslie Lamport.
\textit{\LaTeX: A Document Preparation System}.
Addison-Wesley, 2nd edition, 1994.
\end{thebibliography}

\end{document}
